% =====================================================================
%  Shape morphing: theory, algorithms, and what MorphNode can do about it
%  Build:  pdflatex theory_morphing.tex   (run twice for the ToC)
% =====================================================================
\documentclass[11pt,a4paper]{article}

\usepackage[T1]{fontenc}
\usepackage[utf8]{inputenc}
\usepackage[margin=2.2cm]{geometry}
\usepackage{lmodern}
\usepackage{microtype}
\usepackage{amsmath,amssymb}
\usepackage{booktabs}
\usepackage{enumitem}
\usepackage{xcolor}
\usepackage{hyperref}

\hypersetup{
    colorlinks=true,
    linkcolor=blue!55!black,
    urlcolor=blue!55!black,
    citecolor=blue!55!black,
    pdftitle={Shape morphing for MorphNode},
    pdfauthor={BlenderMathAnim}
}

\newcommand{\src}{S_1}
\newcommand{\dst}{S_2}
\newcommand{\node}[1]{\texttt{#1}}

\title{Shape Morphing in Geometry Nodes\\
       \large Theory, standard algorithms, and a plan for \node{MorphNode}}
\author{BlenderMathAnim --- \node{geometry\_nodes/nodes.py}}
\date{}

\begin{document}
\maketitle
\tableofcontents
\bigskip

% =====================================================================
\section{The problem, concretely}
% =====================================================================

\node{MorphNode} currently computes, for every point $p_i$ of the source
geometry $\src$,
\begin{equation}
  x_i(t) \;=\; (1-t)\,p_i \;+\; t\,q_i ,
  \qquad
  q_i \;=\; \text{position of point $i$ of } \dst ,
  \label{eq:naive}
\end{equation}
with $t$ the \node{Morph Parameter}. Two independent things are wrong with
it, and they are worth separating because they have different fixes.

\paragraph{Vertex counts.}
The pairing $i \mapsto i$ is only defined when
$N = |\src| \le M = |\dst|$. Blender's \node{Sample Index} does \emph{not}
clamp an out-of-range index: it returns the zero vector, so every surplus
source point is sent to the world origin. Measured on a $114$-point UV
sphere morphed into an $8$-point cube: $106$ points land on exactly
$(0,0,0)$. In the picture-frame-into-arrow example of
\node{MorphModifier} the frame carries $4\times 32 = 128$ points and the
arrow $100$, so $28$ points collapse into a knot --- and because the arrow
happens to stand \emph{at} the origin, the knot hides at its foot and reads
as part of the shape.

\paragraph{Edges and faces.}
Worse, and independent of the counts: the index of a point is its position
in a storage array, not a property of the surface. Point $7$ of a swept
rectangle has nothing to do with point $7$ of a cone. Formula
\eqref{eq:naive} therefore transports the source's connectivity along
essentially random trajectories, and the surface folds through itself on the
way. The target's edges and faces never enter the computation at all: what
comes out at $t=1$ is the source's mesh with the target's point positions,
which is only the target if the two meshes were compatible to begin with.

The literature calls these the \emph{correspondence problem} and the
\emph{interpolation} (or \emph{vertex path}) \emph{problem}, and treats them
separately~\cite{alexa2002,lazarus1998}. Everything below follows that split.

% =====================================================================
\section{The correspondence problem}
% =====================================================================

The goal is a bijection $\varphi:\src\to\dst$ between the two surfaces,
together with a mesh that represents both --- a \emph{supermesh} whose
connectivity is valid for the source \emph{and} the target. Given that, the
morph is a per-vertex interpolation and the edges and faces take care of
themselves.

\subsection{Map both shapes to a common domain}

The dominant family. Both surfaces are parameterised over one domain,
correspondence is composition through it, and the supermesh is the overlay
of the two parameterisations.

\begin{itemize}[leftmargin=1.4em]
  \item \textbf{Sphere as domain (genus $0$).} Kent, Carlson and
    Parent~\cite{kent1992} project two star-shaped polyhedra onto a common
    sphere, merge the two projected topologies into one, and interpolate.
    This is the original recipe and still the clearest one to implement.
  \item \textbf{Harmonic and barycentric maps.} For meshes with boundary,
    fix the boundary to a convex polygon and solve a linear system for the
    interior; Tutte's theorem guarantees the embedding is
    valid~\cite{tutte1963}, and Floater's weights give a smooth, shape
    preserving version~\cite{floater1997}. Kanai, Suzuki and Kimura use
    harmonic maps for metamorphosis of arbitrary triangular
    meshes~\cite{kanai2000}.
  \item \textbf{Base domains and consistent parameterisation.} Cut both
    meshes into matching patches over a common coarse base mesh: multi
    resolution analysis of arbitrary meshes~\cite{eck1995}, consistent mesh
    parameterisations~\cite{praun2001}, cross-parameterisation and
    compatible remeshing~\cite{kraevoy2004}, inter-surface
    mapping~\cite{schreiner2004}. These handle higher genus and let the user
    pin feature correspondences (nose to nose, tip to tip).
\end{itemize}

\subsection{Skip the mesh entirely}

If the shapes are represented as functions rather than triangles, no
correspondence is needed: interpolate the \emph{representation}.

\begin{itemize}[leftmargin=1.4em]
  \item \textbf{Implicit functions.} Blend the defining scalar fields and
    re-extract the surface~\cite{turk1999}.
  \item \textbf{Distance fields.} Interpolate signed distance fields, with
    a warp to keep features aligned~\cite{cohenor1998}.
  \item \textbf{Optimal transport.} Treat both shapes as measures and
    follow the Wasserstein geodesic; convolutional Wasserstein distances
    make this practical on geometric domains~\cite{solomon2015}. This is
    the principled answer to ``which point goes where'' when the two point
    sets differ in size, since transport plans are not required to be
    bijections.
\end{itemize}

Topology changes (a torus becoming a ball) come for free here and are
essentially impossible in the mesh-based family.

\subsection{Discrete matching heuristics}

Cheap approximations that do not build a supermesh, which is the class the
node can realistically use:

\begin{itemize}[leftmargin=1.4em]
  \item \textbf{Rigid pre-alignment.} Centre both shapes (and optionally
    normalise scale and principal axes) before matching. This is the
    initialisation step of ICP~\cite{besl1992}; without it every geometric
    matching criterion is dominated by where the two objects happen to sit.
  \item \textbf{Nearest-point projection.} Send each source point to the
    closest point of the target surface. Fast, order-independent, and
    surjective onto nothing in particular: it is many-to-one, so it is a
    shrink-wrap rather than a bijection. It respects geometry, which is
    exactly what index pairing does not.
  \item \textbf{Optimal assignment.} With equal counts, choose the
    permutation minimising $\sum_i \|p_i - q_{\sigma(i)}\|^2$ --- the
    assignment problem, solvable in $O(n^3)$ by the Hungarian
    method~\cite{kuhn1955}. This is the discrete cousin of optimal
    transport and removes the twisting that nearest-point matching leaves
    behind, at the price of a global solve.
\end{itemize}

% =====================================================================
\section{The interpolation problem}
% =====================================================================

Given corresponding points, formula \eqref{eq:naive} moves each along a
straight line. That is the naive choice and it is visibly wrong for large
deformations: a rotating limb shrinks towards the centre of rotation,
because the straight line between two points of a rotating rigid body cuts
the chord instead of following the arc. Standard alternatives:

\begin{itemize}[leftmargin=1.4em]
  \item \textbf{Intrinsic interpolation (2D).} Interpolate edge lengths and
    turning angles instead of coordinates and reconstruct the
    polygon~\cite{sederberg1992,sederberg1993}.
  \item \textbf{As-rigid-as-possible interpolation.} Interpolate the affine
    map of each triangle by splitting it into rotation and stretch, and
    solve a least-squares problem for the vertex positions that best agrees
    with all of them~\cite{alexa2000}. The same machinery underlies ARAP
    surface modelling~\cite{sorkine2007}.
  \item \textbf{Gradient-domain / Poisson blending.} Interpolate the
    deformation gradients and integrate.
\end{itemize}

All three need a global linear solve per frame, which puts them out of
reach of a pure geometry-node graph.

% =====================================================================
\section{What is feasible inside geometry nodes}
% =====================================================================

A geometry-node graph evaluates \emph{fields}: one expression per element,
with no ordering, no iteration to convergence, and no way to build new
connectivity beyond what the fixed operators provide. That rules a lot out.

\begin{center}
\begin{tabular}{@{}lll@{}}
\toprule
\textbf{Method} & \textbf{Needs} & \textbf{In geometry nodes?} \\
\midrule
Index pairing               & nothing                    & yes (today) \\
Index rescaling             & two counts                 & yes \\
Centroid / scale alignment  & attribute statistics       & yes \\
Nearest-point projection    & spatial query              & yes (\node{Geometry Proximity}) \\
Curve arc-length pairing    & resampling                 & yes (curves only) \\
Optimal assignment          & global $O(n^3)$ solve      & no \\
Spherical parameterisation  & iterative relaxation       & no (needs Python) \\
Supermesh construction      & topology surgery           & no (needs Python) \\
ARAP interpolation          & linear solve per frame     & no (needs Python) \\
\bottomrule
\end{tabular}
\end{center}

The three rows marked ``yes'' that are not already there are what the plan
below implements. The rest belong in a Python pre-pass that writes a
compatible pair of meshes once, after which the existing node interpolates
them correctly --- see \S\ref{sec:roadmap}.

% =====================================================================
\section{Implemented plan}
% =====================================================================

Three steps, in the order the graph applies them. Each is a frame in the
node group.

\subsection*{Step 1 --- index rescaling (frame \node{Index Pairing})}

Replace $i \mapsto i$ with the monotone map
\begin{equation}
  j(i) \;=\;
  \operatorname{round}\!\left( i \cdot \frac{M-1}{\,N-1\,} \right),
  \qquad N=|\src|,\; M=|\dst| ,
  \label{eq:rescale}
\end{equation}
clamped to $[0, M-1]$, with $j=0$ when $N \le 1$. The point counts come from
\node{Domain Size} on both inputs, so the map adapts to whatever is plugged
in.

\emph{What it fixes:} the index is never out of range, so nothing is ever
sent to the origin --- the count mismatch stops being a crash-like artefact
and becomes a mere loss of resolution. When $N > M$ several source points
share a destination; when $N < M$ the source samples the target evenly
rather than only its first $N$ points.
\emph{What it does not fix:} the pairing is still storage order.

\subsection*{Step 2 --- centroid alignment (frame \node{Alignment})}

With
\[
  c_k \;=\; \frac{1}{|S_k|}\sum_{p \in S_k} p
  \qquad (\text{\node{Attribute Statistic}, mean of \node{Position}}),
\]
correspondence is computed between the \emph{centred} shapes
$\src - c_1$ and $\dst - c_2$, and the destination is mapped back by $+c_2$.

Note that for a \emph{linear} blend of an \emph{index} pairing this is a
no-op: $(1-t)(p-c_1) + t(q-c_2)$ plus the interpolated centroid
$(1-t)c_1+tc_2$ is exactly $(1-t)p+tq$. It matters only because Step 3
matches by distance, and distance is meaningless while the two shapes sit
metres apart. It is the translation-only case of the rigid pre-registration
that every geometric matcher assumes~\cite{besl1992}.

\subsection*{Step 3 --- nearest-surface pairing (frame \node{Nearest Pairing})}

Pair by geometry instead of by index:
\begin{equation}
  \varphi(p) \;=\;
  c_2 \;+\; \operatorname*{arg\,min}_{y \,\in\, \dst - c_2}
  \bigl\| (p - c_1) - y \bigr\| ,
  \label{eq:proximity}
\end{equation}
evaluated by \node{Geometry Proximity} against the centred target, with
\node{Target Element} $=$ \node{Faces} so the destination may lie anywhere
on the target surface rather than only on its vertices.

\emph{What it fixes:} the destination of a point now depends on where the
point \emph{is}, not on when it was created. Neighbouring source points get
neighbouring destinations, so the source's edges and faces stay coherent
instead of tangling, and the count question disappears entirely --- a
continuous surface has a nearest point for every source point whatever the
two resolutions are.
\emph{What it does not fix:} the map is many-to-one. At $t=1$ the source is
shrink-wrapped onto the target: the silhouette is right and the surface lies
on the target, but concavities the projection cannot see stay uncovered, and
the source's connectivity, not the target's, is what survives.

A \node{Match Nearest} boolean switches between \eqref{eq:rescale} and
\eqref{eq:proximity}; the blend \eqref{eq:naive} is unchanged behind it.

% =====================================================================
\section{The supermesh pass}\label{sec:supermesh}
% =====================================================================

Implemented in \node{geometry\_nodes/supermesh.py}. Everything above is a
compromise forced by the graph; this is the thing that actually solves the
problem, by refusing to pair anything at all.

Both shapes are written over one spherical domain, following Kent, Carlson
and Parent~\cite{kent1992}, and \emph{the domain's own triangulation becomes
the shared connectivity}. Where they merge the two projected topologies into
an exact overlay --- every vertex of both, plus every edge-edge crossing ---
the pass resamples instead: an icosphere is tessellated once at a chosen
subdivision level, and each of its directions $d$ is traced onto both
surfaces,
\begin{equation}
  p_A(d) = c_A + t_A(d)\, d ,
  \qquad
  p_B(d) = c_B + t_B(d)\, d ,
  \label{eq:trace}
\end{equation}
with $t$ found by M\"oller--Trumbore intersection against every triangle,
vectorised and chunked. The \emph{farthest} hit is kept, so a ray leaving a
cone that stands on a cylinder lands on the outside of the pair rather than
on the seam between them.

The trade is explicit. The overlay is exact and its output resolution is a
consequence of the inputs; resampling loses whatever falls between two
directions, but cannot produce degenerate triangles, needs no spherical
polygon clipping, and makes the resolution a parameter --- 642 points at
subdivision 3, 2562 at 4, whatever the inputs carry.

What comes out is two vertex arrays over one face array. Feeding them to
\node{MorphNode} as \node{Geometry 1} and \node{Geometry 2} with
\node{Match Nearest} \emph{off} makes index pairing exactly right, because
point $i$ of the one \emph{is} point $i$ of the other; the node needs no
change. Alternatively one mesh carries the other's positions as a point
attribute, and a \node{Named Attribute} node reads the destination with no
sampling at all.

\paragraph{Assumptions, and how they are reported.}
Equation~\eqref{eq:trace} assumes each shape is star-shaped about its centre
$c$ (the mean of its vertices unless the caller says otherwise) and of genus
$0$. Neither is checked and forbidden; both are measured and reported:

\begin{itemize}[leftmargin=1.4em]
  \item directions that meet no triangle are counted as \emph{misses} and
    fall back to the surface point lying most nearly that way, which keeps
    the map total;
  \item the Euler characteristic $V - E + F$ of each input is computed, since
    $\chi \ne 2$ means no spherical parameterisation exists at all.
\end{itemize}

A torus of revolution run against a cube reports $510$ of $642$ directions
missed and $\chi = 0$, i.e.\ genus~$1$ --- which is the honest answer for the
picture frame of \node{MorphModifier}, and the reason that particular morph
needs a base domain with a handle rather than a
sphere~\cite{praun2001,kraevoy2004,schreiner2004}.

\paragraph{Verified.} A $24\times12$ UV sphere against a $20$-sided cone
gives $642$ points and $1280$ faces with no misses and $\chi = 2$ for both;
the geometry the node evaluates agrees with $(1-t)\,p_A + t\,p_B$ to
$1.2\times10^{-7}$ at $t = 0, \tfrac14, \tfrac12, \tfrac34, 1$, with the face
count constant throughout --- which is what ``the connectivity is the same at
every value of the morph'' means in practice.

% =====================================================================
\section{Roadmap}\label{sec:roadmap}
% =====================================================================

\begin{enumerate}[leftmargin=1.6em]
  \item \textbf{Arc-length pairing for curve-like shapes.} Convert both
    shapes to curves and \node{Resample Curve} to a common count. Both
    problems vanish at once: equal counts, and index order becomes arc
    length, which \emph{is} geometric. Cheap, and enough for outline
    morphs (letters, silhouettes, the picture frame of
    \node{MorphModifier}, which the sphere cannot take).
  \item \textbf{Higher genus and feature control.} Replace the sphere of
    \S\ref{sec:supermesh} by a common base domain with the right handles,
    and let the caller pin correspondences (tip to tip)
    ~\cite{praun2001,kraevoy2004,schreiner2004}.
  \item \textbf{Better trajectories.} With the supermesh in place, replace
    the straight line by an as-rigid-as-possible path~\cite{alexa2000}; a
    per-frame least-squares solve, so again Python.
\end{enumerate}

% =====================================================================
\begin{thebibliography}{99}
\small

\bibitem{alexa2000} M.~Alexa, D.~Cohen-Or, D.~Levin.
  \emph{As-rigid-as-possible shape interpolation}. SIGGRAPH, 2000.

\bibitem{alexa2002} M.~Alexa.
  \emph{Recent advances in mesh morphing}. Computer Graphics Forum 21(2), 2002.

\bibitem{besl1992} P.~J.~Besl, N.~D.~McKay.
  \emph{A method for registration of 3-D shapes}. IEEE Transactions on
  Pattern Analysis and Machine Intelligence 14(2), 1992.

\bibitem{cohenor1998} D.~Cohen-Or, D.~Levin, A.~Solomovici.
  \emph{Three-dimensional distance field metamorphosis}. ACM Transactions on
  Graphics 17(2), 1998.

\bibitem{eck1995} M.~Eck, T.~DeRose, T.~Duchamp, H.~Hoppe, M.~Lounsbery,
  W.~Stuetzle. \emph{Multiresolution analysis of arbitrary meshes}.
  SIGGRAPH, 1995.

\bibitem{floater1997} M.~S.~Floater.
  \emph{Parametrization and smooth approximation of surface triangulations}.
  Computer Aided Geometric Design 14(3), 1997.

\bibitem{kanai2000} T.~Kanai, H.~Suzuki, F.~Kimura.
  \emph{Metamorphosis of arbitrary triangular meshes}. IEEE Computer
  Graphics and Applications 20(2), 2000.

\bibitem{kent1992} J.~R.~Kent, W.~E.~Carlson, R.~E.~Parent.
  \emph{Shape transformation for polyhedral objects}. SIGGRAPH, 1992.

\bibitem{kraevoy2004} V.~Kraevoy, A.~Sheffer.
  \emph{Cross-parameterization and compatible remeshing of 3D models}.
  ACM Transactions on Graphics (SIGGRAPH), 2004.

\bibitem{kuhn1955} H.~W.~Kuhn.
  \emph{The Hungarian method for the assignment problem}. Naval Research
  Logistics Quarterly 2, 1955.

\bibitem{lazarus1998} F.~Lazarus, A.~Verroust.
  \emph{Three-dimensional metamorphosis: a survey}. The Visual Computer
  14(8--9), 1998.

\bibitem{praun2001} E.~Praun, W.~Sweldens, P.~Schr\"oder.
  \emph{Consistent mesh parameterizations}. SIGGRAPH, 2001.

\bibitem{schreiner2004} J.~Schreiner, A.~Asirvatham, E.~Praun, H.~Hoppe.
  \emph{Inter-surface mapping}. ACM Transactions on Graphics (SIGGRAPH), 2004.

\bibitem{sederberg1992} T.~W.~Sederberg, E.~Greenwood.
  \emph{A physically based approach to 2-D shape blending}. SIGGRAPH, 1992.

\bibitem{sederberg1993} T.~W.~Sederberg, P.~Gao, G.~Wang, H.~Mu.
  \emph{2-D shape blending: an intrinsic solution to the vertex path
  problem}. SIGGRAPH, 1993.

\bibitem{solomon2015} J.~Solomon, F.~de~Goes, G.~Peyr\'e, M.~Cuturi,
  A.~Butscher, A.~Nguyen, T.~Du, L.~Guibas.
  \emph{Convolutional Wasserstein distances: efficient optimal
  transportation on geometric domains}. ACM Transactions on Graphics
  (SIGGRAPH), 2015.

\bibitem{sorkine2007} O.~Sorkine, M.~Alexa.
  \emph{As-rigid-as-possible surface modeling}. Symposium on Geometry
  Processing, 2007.

\bibitem{turk1999} G.~Turk, J.~F.~O'Brien.
  \emph{Shape transformation using variational implicit functions}.
  SIGGRAPH, 1999.

\bibitem{tutte1963} W.~T.~Tutte.
  \emph{How to draw a graph}. Proceedings of the London Mathematical
  Society, 1963.

\end{thebibliography}

\end{document}